\chapter{Annexes}
\thispagestyle{empty}

\section{Gains relatifs Eval 3.2}

Le Tableau~\ref{tab:annex-relative-gains} resume l'analyse des gains relatifs en
pourcentage introduite dans l'extension Eval 3.2. Les valeurs positives
indiquent une amelioration dans le sens approprie de la metrique.

\begin{table}[H]
\centering
\caption{Ameliorations relatives (\%) issues de l'extension Eval 3.2.}
\label{tab:annex-relative-gains}
\setlength{\tabcolsep}{4pt}
\renewcommand{\arraystretch}{1.12}
\resizebox{\textwidth}{!}{
\begin{tabular}{@{}l*{11}{c}@{}}
\toprule
\textbf{Comparaison} &
\textbf{H@10} &
\textbf{NDCG@10} &
\textbf{Recall@10} &
\textbf{MAP@10} &
\textbf{ROC-AUC} &
\textbf{PR-AUC} &
\textbf{Pop. Lift} &
\textbf{Gini} &
\textbf{Agg. Div.} &
\textbf{Novelty} &
\textbf{ILD} \\
\midrule
Hybride principal vs texte seul & 14.60 & 0.77 & 13.55 & -8.46 & 8.04 & 7.58 & 31.48 & 20.53 & 3.86 & 3.09 & 17.75 \\
Hybride principal vs graphe seul & 11.03 & 15.22 & 11.02 & 19.26 & -1.58 & -1.60 & 0.54 & 8.29 & 3.28 & 1.62 & -4.30 \\
Effet du pooling (corrige) & 0.05 & -0.95 & 0.05 & -1.81 & -0.19 & -0.22 & -1.68 & 0.41 & 0.51 & -0.03 & 1.31 \\
Effet du pooling (non corrige) & 0.01 & -0.04 & 0.01 & -0.09 & -0.00 & -0.00 & 0.02 & -0.17 & 0.20 & 0.01 & 0.03 \\
Effet de la correction (Node2Vec) & -0.04 & -0.12 & -0.04 & -0.21 & -0.00 & -0.00 & 0.00 & 0.00 & 0.00 & 0.00 & 0.00 \\
Effet de la correction (hybride) & 10.66 & 16.00 & 10.66 & 21.12 & -1.32 & -1.40 & 2.38 & 10.38 & 0.93 & 1.74 & -5.30 \\
Effet de la correction (hybride poolé) & 10.70 & 14.95 & 10.70 & 19.04 & -1.51 & -1.61 & 0.71 & 10.90 & 1.25 & 1.70 & -4.08 \\
\bottomrule
\end{tabular}
}
\end{table}

\section{Contraintes de l'organisateur LLM}

L'organisateur LLM n'est pas autorise a produire librement un nouveau contenu
de recommandation. Dans l'application finale, il doit respecter un contrat de
sortie strict, qui peut etre resume ainsi:
\begin{itemize}
    \item utiliser uniquement les 20 titres deja fournis par le moteur de
    recommandation;
    \item affecter chaque titre a une seule section parmi
    \textit{Foundations}, \textit{Core Concepts}, \textit{Advanced Topics} et
    \textit{Adjacent Topics};
    \item conserver chaque rang exactement une fois et ne jamais dupliquer un
    article;
    \item fournir pour chaque affectation une justification courte compatible
    avec la section choisie.
\end{itemize}

Si ce contrat n'est pas respecte, la sortie est rejetee et l'application
emploie le repli heuristique decrit dans le corps du rapport. Ce choix permet
de conserver un comportement reproductible et une presentation exploitable sans
faire du LLM le coeur de la recommandation.
